\documentclass[11pt]{article}

\usepackage[T1]{fontenc}
\usepackage[utf8]{inputenc}
\usepackage[english]{babel}
\usepackage{amsmath,amssymb,mathtools,amsthm}
\usepackage{enumitem}
\usepackage[a4paper,margin=2.5cm]{geometry}
\usepackage[hidelinks]{hyperref}
\usepackage{microtype}

% -----------------------------------------------------------------------------
% Notation
% -----------------------------------------------------------------------------
\newcommand{\C}{\mathbb{C}}
\newcommand{\N}{\mathbb{N}}
\newcommand{\BH}{\mathcal{B}(H)}
\newcommand{\FH}{\mathcal{F}(H)}
\newcommand{\KH}{\mathcal{K}(H)}
\newcommand{\QH}{\mathcal{Q}(H)}
\newcommand{\spann}{\operatorname{span}}
\newcommand{\rank}{\operatorname{rank}}
\newcommand{\ran}{\operatorname{ran}}

% Simple environments: deliberately without boxes or colored backgrounds.
\newenvironment{definition}{\par\medskip\noindent\textbf{Definition.}\ }{\par\medskip}
\newenvironment{proposition}{\par\medskip\noindent\textbf{Proposition.}\ }{\par\medskip}
\newenvironment{observation}{\par\medskip\noindent\textbf{Observation.}\ }{\par\medskip}
\newenvironment{exercise}{\par\medskip\noindent\textbf{Exercise.}\ }{\par\medskip}
\newenvironment{calculation}{\par\medskip\noindent\textbf{Calculation.}\ }{\par\medskip}

\title{Continuation of the $C^*$-Algebra Notes\\Finite-Dimensional Corners, Compact Operators, and the Calkin Algebra}
\author{}
\date{}

\begin{document}
\maketitle

\noindent
These notes continue the rank-one and corner calculations developed in the previous
session.  The handwritten order has been preserved, while a few short explanations
have been added to make the transitions explicit.  Such additions are marked in the
text or in footnotes.

Throughout, let $H$ be an infinite-dimensional separable complex Hilbert space with
fixed orthonormal basis
\[
(e_1,e_2,\ldots).
\]
We keep the convention from the preceding notes that the inner product is linear in
the first variable, and
\[
(e_i\otimes e_j)(h)=\langle h,e_j\rangle e_i.
\]
The assumption that the basis is indexed by $\N$ is exactly the separability assumption
used implicitly in the handwritten notes.\footnote{For a nonseparable Hilbert space one
usually replaces the sequence $(p_n)$ below by a directed family of finite-dimensional
projections.  The definition of $\KH$ remains the same, but some statements about the
ideal structure of $\BH$ require additional care.}

\section{The rank-two corner revisited}

Let
\[
p=e_1\otimes e_1+e_2\otimes e_2.
\]
If $S\in\BH$, write
\[
S(e_j)=\sum_{k=1}^{\infty}\alpha_{kj}e_k,
\qquad j=1,2.
\]
Using the matrix-unit relation
\[
(e_i\otimes e_j)(e_r\otimes e_s)
 =\delta_{jr}\,e_i\otimes e_s,
\]
one obtains
\begin{align*}
pSp
&=\alpha_{11}e_1\otimes e_1
 +\alpha_{12}e_1\otimes e_2
 +\alpha_{21}e_2\otimes e_1
 +\alpha_{22}e_2\otimes e_2.
\end{align*}
Therefore the map
\[
pSp\longmapsto
\begin{pmatrix}
\alpha_{11}&\alpha_{12}\\
\alpha_{21}&\alpha_{22}
\end{pmatrix}
\]
identifies the corner $p\BH p$ with $M_2(\C)$ as a $C^*$-algebra.\footnote{In the
handwritten notes the four coefficients are also displayed as a point of $\C^4$.
That is correct as a \emph{vector-space} identification.  The algebra product,
however, is the matrix product, so the natural algebraic identification is
$p\BH p\cong M_2(\C)$, not coordinatewise multiplication on $\C^4$.}

\section{The finite-dimensional projections $p_n$}

For $n\in\N$, define
\[
p_n:=\sum_{i=1}^{n}e_i\otimes e_i.
\]
Then $p_n$ is the orthogonal projection onto
\[
H_n:=\spann\{e_1,\ldots,e_n\}.
\]

\begin{proposition}
For every $n\ge1$,
\[
p_n\BH p_n\cong M_n(\C).
\]
\end{proposition}

\noindent
Indeed, if
\[
T(e_j)=\sum_{k=1}^{\infty}\alpha_{kj}e_k,
\]
then
\[
p_nTp_n
 =\sum_{i,j=1}^{n}\alpha_{ij}\,e_i\otimes e_j.
\]
Hence the $*$-isomorphism is given by
\[
\Phi_n:p_n\BH p_n\longrightarrow M_n(\C),
\qquad
\Phi_n(p_nTp_n)=[\alpha_{ij}]_{i,j=1}^{n}.
\]
The unit of the corner is $p_n$, and under $\Phi_n$ it corresponds to $I_n$.
This is the standard realization of $M_n(\C)$ as a finite-dimensional corner of
$\BH$.\footnote{Compare Davidson, \cite[\S I.10]{Davidson1996}, and Erdman,
\cite[\S 7.3--7.4]{Erdman2011}.}

\subsection{Compatibility of the corners}

The handwritten notes record the inclusion
\[
p_n\BH p_n\subseteq p_{n+1}\BH p_{n+1}.
\]
To see it, first note that
\begin{align*}
p_np_{n+1}
&=\left(\sum_{i=1}^{n}e_i\otimes e_i\right)
  \left(\sum_{j=1}^{n+1}e_j\otimes e_j\right)\\
&=\sum_{i=1}^{n}e_i\otimes e_i
 =p_n,
\end{align*}
and similarly $p_{n+1}p_n=p_n$.
Thus, for $T\in\BH$,
\begin{align*}
p_{n+1}(p_nTp_n)p_{n+1}
 &=p_nTp_n.
\end{align*}
Therefore $p_nTp_n$ belongs to the $(n+1)$-st corner.

Under the identifications $p_n\BH p_n\cong M_n(\C)$, this is exactly the usual
block embedding
\[
M_n(\C)\hookrightarrow M_{n+1}(\C),
\qquad
x\longmapsto
\begin{pmatrix}
x&0\\
0&0
\end{pmatrix}.
\]
Consequently,
\[
\bigcup_{n=1}^{\infty}M_n(\C)
\cong
\bigcup_{n=1}^{\infty}p_n\BH p_n
\subseteq \BH.
\]
This union is an increasing $*$-subalgebra of $\BH$ consisting of operators whose
matrix, relative to the chosen basis, has support in a finite upper-left block.\footnote{For
a fixed basis this algebraic union need not equal all of $\FH$: a rank-one operator
$\xi\otimes\eta$ may involve vectors $\xi,\eta$ with infinitely many nonzero basis
coefficients.  What is true, and what is used below, is that the norm closure of this
increasing union is $\KH$.}

\section{Finite-rank operators}

Let $\FH$ denote the linear space of finite-rank operators in $\BH$.

\begin{proposition}
If $T\in\FH$, then $T$ is a finite sum of rank-one operators.  More precisely,
if
\[
\{e_1,\ldots,e_n\}
\]
is an orthonormal basis for $\ran T$, then
\[
\boxed{
T=\sum_{i=1}^{n}e_i\otimes T^*(e_i).
}
\]
\end{proposition}

\begin{proof}
For $h\in H$, since $Th\in\ran T$,
\begin{align*}
Th
 &=\sum_{i=1}^{n}\langle Th,e_i\rangle e_i\\
 &=\sum_{i=1}^{n}\langle h,T^*e_i\rangle e_i\\
 &=\left(\sum_{i=1}^{n}e_i\otimes T^*(e_i)\right)(h).
\end{align*}
Therefore
\[
T=\sum_{i=1}^{n}e_i\otimes T^*(e_i).
\]
\end{proof}

This proves, in particular, that the finite-rank operators are exactly the finite
linear combinations of rank-one operators.\footnote{A standard reference is Conway,
\cite[Chapter II, \S4]{Conway1990}; see also the explicit rank-one discussion in
Erdman, \cite[\S 7.3]{Erdman2011}.}

\section{The algebra of compact operators}

\begin{definition}
The $C^*$-algebra of compact operators on $H$ is
\[
\boxed{
\KH:=\overline{\FH}^{\,\|\cdot\|}.
}
\]
That is, $\KH$ is the operator-norm closure of the finite-rank operators.
\end{definition}

On a Hilbert space this norm-closure definition agrees with the usual topological
definition: $T\in\BH$ is compact if it maps the closed unit ball into a relatively
compact subset of $H$.\footnote{This equivalence is one of the standard basic results
on compact operators; see Conway, \cite[Chapter II, Theorem 4.4]{Conway1990}, or
MIT 18.102, Lecture 18, \cite{MIT2009}.}

\subsection{Compact operators from the matrix corners}

The finite-dimensional corners form a norm-dense subalgebra of $\KH$:
\[
\boxed{
\KH
 =\overline{\bigcup_{n=1}^{\infty}p_n\BH p_n}^{\,\|\cdot\|}
 \cong
 \overline{\bigcup_{n=1}^{\infty}M_n(\C)}^{\,\|\cdot\|}.
}
\]
Here each $M_n(\C)$ is embedded in $M_{n+1}(\C)$ by
\[
x\mapsto
\begin{pmatrix}
x&0\\
0&0
\end{pmatrix}.
\]
The closure is essential: the algebraic union contains only finite matrices, whereas
$\KH$ also contains norm limits of such matrices.\footnote{This viewpoint is the
first model of an inductive-limit construction and is closely related to the later
theory of AF $C^*$-algebras.  Davidson's text develops this point of view systematically;
see \cite[Chapters I and III]{Davidson1996}.}

A useful approximation statement is
\[
T\in\KH
\quad\Longrightarrow\quad
\|T-p_nTp_n\|\longrightarrow0.
\]
For a compact operator this follows from the strong convergence $p_n\to I$ together
with uniform convergence on compact subsets.  This supplies the link between the
handwritten increasing matrix corners and the full compact-operator algebra.

\subsection{$*$-closedness}

Suppose $T_m\in\FH$ and
\[
T_m\longrightarrow T
\qquad\text{in operator norm}.
\]
Since finite rank is preserved by taking adjoints,
\[
T_m^*\in\FH.
\]
Moreover,
\[
\|T_m^*-T^*\|=\|T_m-T\|\longrightarrow0.
\]
Therefore
\[
T\in\KH\quad\Longrightarrow\quad T^*\in\KH.
\]
Thus $\KH$ is closed under the involution.

The inclusions
\[
M_n(\C)\subseteq\KH
\qquad(n\ge1)
\]
show that $\KH$ is infinite-dimensional whenever $H$ is infinite-dimensional, since
\[
\dim_{\C}M_n(\C)=n^2\longrightarrow\infty.
\]

\section{$\KH$ as an ideal in $\BH$}

\begin{proposition}
If $T\in\KH$ and $S\in\BH$, then
\[
TS\in\KH
\qquad\text{and}\qquad
ST\in\KH.
\]
Consequently, $\KH$ is a closed two-sided $*$-ideal of $\BH$.
\end{proposition}

\begin{proof}
Choose finite-rank operators $T_m$ such that
\[
\|T_m-T\|\longrightarrow0.
\]
Then $T_mS$ and $ST_m$ are finite rank, while
\begin{align*}
\|T_mS-TS\|
 &\le \|T_m-T\|\,\|S\|\longrightarrow0,\\
\|ST_m-ST\|
 &\le \|S\|\,\|T_m-T\|\longrightarrow0.
\end{align*}
Hence $TS,ST\in\KH$.
\end{proof}

The handwritten sentence that $\KH$ is ``the unique $C^*$-ideal of $\BH$'' should be
read more precisely as follows.\footnote{The ideals $\{0\}$ and $\BH$ are always present,
so ``unique'' here means unique \emph{nonzero proper norm-closed two-sided ideal}.}

\begin{proposition}
If $H$ is separable and infinite-dimensional, then $\KH$ is the unique nonzero proper
norm-closed two-sided ideal of $\BH$.
\end{proposition}

This separability hypothesis matters: for nonseparable Hilbert spaces, $\BH$ has
additional nontrivial closed ideals.\footnote{For a precise statement, see Blackadar,
\cite[\S I.8.7]{Blackadar2006}.  In particular, he explicitly distinguishes the
separable and nonseparable cases.}

\section{The Calkin algebra}

\begin{definition}
The \emph{Calkin algebra} is the quotient $C^*$-algebra
\[
\boxed{
\QH:=\BH/\KH.
}
\]
If
\[
\pi:\BH\longrightarrow\QH,
\qquad
\pi(T)=T+\KH,
\]
denotes the quotient map, then the inclusion
\[
\iota:\KH\hookrightarrow\BH
\]
and $\pi$ give the short exact sequence
\[
\boxed{
0\longrightarrow\KH
\xrightarrow{\ \iota\ }
\BH
\xrightarrow{\ \pi\ }
\QH
\longrightarrow0.
}
\]
\end{definition}

Thus the Calkin algebra identifies two bounded operators whenever they differ by a
compact operator.  In this sense it records operator-theoretic information ``modulo
compact perturbations.''\footnote{For later study: Atkinson's theorem says that
$T\in\BH$ is Fredholm if and only if its coset $\pi(T)$ is invertible in the Calkin
algebra.  This theorem was not written in the photographs, but it explains why the
quotient $\BH/\KH$ becomes central in Fredholm theory; see Murphy,
\cite[\S 1.4]{Murphy1990}, or Blackadar, \cite[Chapter I]{Blackadar2006}.}

\section*{Reading guide for the photographs}

The mathematical chain represented in the photographs is
\[
\boxed{
\begin{gathered}
 p_n\BH p_n\cong M_n(\C),
 \qquad
 M_n(\C)\hookrightarrow M_{n+1}(\C),\\[2mm]
 \bigcup_{n\ge1}p_n\BH p_n
 \ \text{is norm dense in}\ \KH,\\[2mm]
 \FH
 =\spann\{\text{rank-one operators}\},
 \qquad
 \KH=\overline{\FH}^{\|\cdot\|},\\[2mm]
 \KH\triangleleft\BH,
 \qquad
 \QH=\BH/\KH,\\[2mm]
 0\longrightarrow\KH\longrightarrow\BH
 \longrightarrow\QH\longrightarrow0.
\end{gathered}
}
\]
The sequence should be read as a passage from finite matrices to finite-rank
operators, then to their norm closure, and finally to the quotient that forgets compact
perturbations.

\section*{References and where to look}

\begin{thebibliography}{9}

\bibitem{Conway1990}
J.~B. Conway,
\emph{A Course in Functional Analysis}, 2nd ed.,
Graduate Texts in Mathematics 96, Springer, 1990.
See Chapter II, especially \S4 ``Compact Operators.''

\bibitem{Murphy1990}
G.~J. Murphy,
\emph{$C^*$-Algebras and Operator Theory},
Academic Press, 1990.
See \S1.4 for compact and Fredholm operators and Chapter 2 for Hilbert-space
operator $C^*$-algebras.

\bibitem{Davidson1996}
K.~R. Davidson,
\emph{$C^*$-Algebras by Example},
Fields Institute Monographs 6, American Mathematical Society, 1996.
See \S I.10, ``$C^*$-Algebras of Compact Operators,'' and Chapter III for AF
algebras and inductive-limit ideas.

\bibitem{Blackadar2006}
B.~Blackadar,
\emph{Operator Algebras: Theory of $C^*$-Algebras and von Neumann Algebras},
Encyclopaedia of Mathematical Sciences 122, Springer, 2006.
See \S I.8 on compact operators and \S I.8.7 on ideals of $\mathcal B(H)$.

\bibitem{Erdman2011}
J.~M. Erdman,
\emph{Lecture Notes on Operator Algebras},
Portland State University, 2011.
See \S7.3 ``Finite Rank Operators'' and \S7.4 ``Compact Operators.''
Available online from Portland State University.

\bibitem{MIT2009}
R.~B. Melrose,
\emph{18.102 Introduction to Functional Analysis --- Lecture 18: Compact Operators},
MIT OpenCourseWare, Spring 2009.
The lecture proves that compact operators on a separable Hilbert space are exactly
the norm limits of finite-rank operators.

\end{thebibliography}

\end{document}
